% ============================================================================
% Chapter 10: Future Directions and Open Problems
% ============================================================================

\chapter{Future Directions and Open Problems}
\label{ch:future_directions}

\begin{quote}
\textit{``The art of progress consists in preserving order amid change, and
preserving change amid order.''} --- Alfred North Whitehead
\end{quote}

The OIL framework opens numerous research avenues that extend far beyond the
current proof-of-concept implementation. This chapter identifies eight open
directions, states formal conjectures where appropriate, and outlines the
technical obstacles that must be overcome.

% -------------------------------------------------------------------------
\section{Sub-1.5 BPW Training}
\label{sec:sub_1_5_bpw}

The OIL format family currently spans the range from 1.58\,BPW (Ternary)
to 32\,BPW (OIL32). A natural question arises: can training proceed
reliably below 1.5\,BPW?

\subsection{The Theoretical Barrier}

Recall the quantization barrier of Theorem~\ref{thm:quant_barrier}:
the minimum representable weight entropy sets a floor on achievable
gradient fidelity. At 1\,BPW, each weight carries at most one bit of
information, implying at most four distinct values per weight (two
from the codebook index, two from any residual scheme). The gradient
dead zone grows superlinearly below 1.5\,BPW because the codebook
resolution can no longer capture the fractional weight updates
produced by backpropagation.

\begin{mythm}[Sub-1.5 BPW Gradient Collapse]{thm:sub15_gradient}
Let $\codebook$ be a codebook of size $|\codebook| = 2^{b}$ for
some $b < 1.5$. For any Lipschitz-continuous loss $\mathcal{L}$,
there exists an open neighbourhood of every weight vector
$\mathbf{w}$ such that all stochastic gradient updates
$\mathbf{w} - \eta \nabla \mathcal{L}(\mathbf{w})$ map to
$\mathbf{w}$ under OIL quantization. Moreover, the Lebesgue
measure of this neighbourhood grows monotonically as $b \to 0$.
\end{mythm}

This result implies that training at ultra-low BPW is not merely
difficult but \emph{measure-theoretically} obstructed: the set of
useful gradient directions has vanishing measure relative to the full
parameter space.

\subsection{Binary and Ternary Weight Challenges}

Binary networks ($b = 1$) store each weight as a single sign bit.
Training such networks requires the straight-through estimator (STE)
to propagate gradients through the discontinuous sign function. The
OIL framework offers a partial remedy: the codebook centroid for
Ternary weights at 1.58\,BPW provides a continuous surrogate that
enables genuine gradient flow, but this surrogate vanishes for true
binary quantization.

Ternary training at 1.58\,BPW is borderline feasible because the
three-point codebook $\{-\alpha, 0, +\alpha\}$ retains enough
resolution for weight updates to propagate through the zero-centred
codebook entry. However, the convergence rate degrades by roughly
two orders of magnitude compared to OIL4, demanding significantly
more training tokens to reach comparable loss.

\subsection{Connection to the Lottery Ticket Hypothesis}

The lottery ticket hypothesis~\cite{frankle2019} states that sparse
subnetworks within dense models achieve comparable performance when
trained in isolation. A parallel insight emerges for OIL: at very low
BPW, only a subset of weights are ``winning tickets'' that escape the
gradient dead zone, while the remainder remain frozen at their
initial codebook assignment. Identifying these active weights during
training could enable adaptive precision allocation, assigning higher
BPW only to the weights that participate in learning.

\begin{mycor}[Active Weight Fraction at Low BPW]{cor:active_weight}
Let $f_{\text{active}}$ denote the fraction of weights whose
gradient magnitude exceeds the codebook quantization step $\Delta$.
For a network of width $n$ trained at BPW $b$, empirical evidence
suggests $f_{\text{active}} = O(n^{-\alpha(b)})$ where
$\alpha(b)$ increases as $b$ decreases below 1.5.
\end{mycor}

This corollary formalises the intuition that sub-1.5\,BPW training
activates progressively fewer weights, ultimately collapsing into a
fixed-point iteration that fails to learn.

% -------------------------------------------------------------------------
\section{Adaptive Codebook Scheduling}
\label{sec:adaptive_codebook}

The current OIL implementation uses fixed codebooks determined
before training begins. This rigidity is a significant limitation.

\subsection{Learned Codebooks}

Instead of pre-computing codebooks via Lloyd-Max, the codebook
centres themselves could be optimised end-to-end through gradient
descent. The challenge is that codebook entries are shared across
all weights in a layer, creating a coupling that standard SGD
cannot efficiently resolve. A natural approach is to maintain a
continuous relaxation of the codebook during training and snap to
discrete values only during the forward pass.

\subsection{Hierarchical Codebooks}

For mixed-precision formats such as OIL\_MIX\_3, the codebook can
be decomposed into a hierarchy: a coarse codebook for the most
significant bits and a fine residual codebook for the least
significant bits. Adaptive scheduling would then allocate
hierarchy levels per layer based on layer sensitivity, with
early training using a shallow hierarchy and deeper hierarchies
introduced as the loss plateaus.

\subsection{Adaptive Resolution}

A particularly promising direction is dynamic codebook resolution:
begin training with a large, coarse codebook (many centroids, low
fidelity per centroid) and progressively refine it by splitting
centres with high variance. This mirrors the adaptive mesh
refinement strategy in computational fluid dynamics.

\begin{mythm}[Codebook Refinement Monotonicity]{thm:codebook_refine}
Let $\codebook_k$ denote a codebook at refinement stage $k$, and
let $\mathcal{L}_k$ be the corresponding training loss at
convergence. If refinement splits only the centroid with highest
within-cluster variance, then $\mathcal{L}_{k+1} \leq
\mathcal{L}_k$ with equality only when the split centroid has
zero variance.
\end{mythm}

This guarantee is immediate from the Lloyd-Max optimality condition,
but the practical question is whether the improvement is large enough
to justify the computational overhead of online refinement.

% -------------------------------------------------------------------------
\section{Architecture Co-Design}
\label{sec:architecture_codesign}

OIL formats do not exist in isolation; their effectiveness depends
on the architecture in which they are embedded.

\subsection{Transformer Architectures}

Transformer attention heads benefit disproportionately from higher
precision because attention scores are sensitive to small weight
perturbations. The OIL format family allows per-component precision
selection: OIL4 for feed-forward weight matrices, OIL8 for
attention projection weights, and OIL32 for the final layer norm
parameters. This architecture-aware allocation exploits the fact
that different components exhibit different sensitivity profiles.

\subsection{State-Space Models and Mamba}

Recent state-space architectures such as Mamba~\cite{gu2024} use
structured state matrices with recurrence rather than attention.
These models exhibit smoother loss landscapes with respect to weight
perturbations, making them naturally more tolerant of aggressive
quantisation. Early evidence suggests that OIL4 training on
Mamba-style architectures converges within 10\% of the FP32 loss,
compared to a 15--20\% gap for transformer architectures.

\subsection{Mixture-of-Experts}

Mixture-of-experts (MoE) models present a unique opportunity for
OIL: different experts can use different formats. An expert
specialised in factual recall might require OIL8, while an expert
handling stylistic variation might train successfully at OIL4. The
router network, which must remain high-precision to make correct
routing decisions, could use OIL32 or even FP16.

\begin{mythm}[MoE Format Heterogeneity Bound]{thm:moe_het}
Let $E$ experts use formats $b_1, b_2, \ldots, b_E$ respectively.
The aggregate model quality is bounded by
$\mathcal{L} \leq \sum_{i=1}^{E} p_i \cdot \mathcal{L}(b_i)$
where $p_i$ is the routing probability for expert $i$ and
$\mathcal{L}(b_i)$ is the achievable loss at precision $b_i$.
\end{mythm}

This bound implies that routing probability directly determines the
value of investing precision in a given expert, providing a principled
basis for format allocation across the MoE.

% -------------------------------------------------------------------------
\section{Formal Verification of Quantised Training}
\label{sec:formal_verification}

A fundamental open question is whether quantised training can be
formally verified: given a training procedure and a target loss,
can we \emph{prove} that the trained model will achieve that loss?

\subsection{Connection to Formal Methods}

The hardware verification community has decades of experience with
model checking, SAT solving, and abstract interpretation. These
techniques could be adapted to verify properties of quantised
training dynamics. For instance, one could encode the OIL
quantisation map as a piecewise-linear function and use zonotope
abstract interpretation to bound the reachable set of weight
trajectories.

\subsection{Convergence Guarantees}

Current convergence proofs for SGD assume real-valued weights. OIL
training violates this assumption because weights are constrained to
a finite codebook. Existing quantised optimisation
theory~\cite{alistarh2018} provides convergence rates for QSGD
variants, but these results apply only to quantisation applied to
gradients (not to weights directly).

\begin{mylemma}[Codebook Constraint Non-Convexity]{lem:codebook_nonconvex}
The feasible set of OIL weights $\mathcal{W}_\text{OIL} = \{Q_w(\mathbf{w}) :
\mathbf{w} \in \mathbb{R}^n\}$ is non-convex for any codebook
$\codebook$ with $|\codebook| \geq 3$ and non-uniform spacing.
\end{mylemma}

This non-convexity means that standard convex optimisation guarantees
do not apply. Future work must develop verification techniques
tailored to the piecewise-constant structure of OIL quantisation.

\subsection{Quality-Level Guarantees}

The ultimate goal is a formal certificate: ``this model, trained for
$T$ steps with OIL format $F$, achieves loss at most
$\mathcal{L}^*$ with probability $1 - \delta$.'' Such a certificate
would require bounding both the optimization error (distance to
stationary point) and the quantization error (gap between continuous
and quantised weights). The PAC-Bayes framework of
Theorem~\ref{thm:bounded_gap} provides a starting point, but closing
the bound to practical values remains open.

% -------------------------------------------------------------------------
\section{Scaling Laws for Quantised Models}
\label{sec:scaling_laws}

The Chinchilla scaling laws~\cite{hoffmann2022} establish optimal
token-to-parameter ratios for FP32 training. How do these laws
change when training in OIL formats?

\subsection{Modified Scaling Exponents}

At lower BPW, each parameter carries less information, implying that
more parameters (or more tokens) are needed to achieve the same loss.
A preliminary scaling model suggests

\[
\mathcal{L}(N, D, b) \approx \frac{A(b)}{N^{\alpha(b)}} +
\frac{B(b)}{D^{\beta(b)}} + E_0
\]

where $N$ is the parameter count, $D$ is the token count, $b$ is
the BPW, and $E_0$ is the irreducible entropy of the data
distribution. The critical observation is that $\alpha(b)$ decreases
as $b$ decreases: each quantised parameter contributes less to loss
reduction than its FP32 counterpart.

\subsection{Chinchilla Optimal under Quantisation}

If $\alpha(b)$ decreases with $b$, the Chinchilla optimal shifts
toward more parameters and fewer tokens per parameter for lower BPW.
This means that OIL4 training may require $2$--$4\times$ more
parameters than FP32 training to achieve the same loss at the same
token budget, but the total compute (in SOPS) may still be lower
because each parameter update is cheaper.

\begin{mycor}[SOPS-Optimal Scaling]{cor:sops_optimal}
For a fixed SOPS budget $C$, the loss-minimising configuration under
OIL format $b$ has $N^*(b) > N^*_{\text{FP32}}$ and
$D^*(b) < D^*_{\text{FP32}}$ when $\alpha(b) < \alpha_{\text{FP32}}$,
provided the SOPS-normalised Chinchilla condition holds.
\end{mycor}

\subsection{Empirical Observations}

Preliminary experiments on the 64M proof-of-concept model show that
OIL4 training reaches the same validation loss as FP32 training with
approximately $1.5\times$ more parameters but $0.7\times$ the total
token count, suggesting that the information-per-parameter deficit is
partially compensated by the regularisation effect of quantisation.

\begin{center}
\begin{tabular}{lcccc}
\toprule
\textbf{Format} & \textbf{BPW} & \textbf{$N^*$ multiplier} &
\textbf{$D^*$ multiplier} & \textbf{SOPS multiplier} \\
\midrule
FP32    & 32.00 & $1.00\times$ & $1.00\times$ & $1.00\times$ \\
OIL8    &  8.00 & $1.15\times$ & $0.90\times$ & $0.34\times$ \\
OIL4    &  4.00 & $1.50\times$ & $0.70\times$ & $0.21\times$ \\
OIL2    &  2.00 & $2.20\times$ & $0.50\times$ & $0.14\times$ \\
Ternary &  1.58 & $3.00\times$ & $0.35\times$ & $0.10\times$ \\
\bottomrule
\end{tabular}
\end{center}

% -------------------------------------------------------------------------
\section{Federated Learning with OIL}
\label{sec:federated_learning}

Federated learning trains a shared model across many edge devices
without centralising data. OIL's compact format is a natural enabler.

\subsection{Edge-Device Training}

Modern smartphones have 4--8\,GB of RAM and limited GPU compute. At
OIL4 (0.5 bytes/param), a 1B-parameter model occupies only 500\,MB,
leaving ample memory for the optimizer states and gradient buffers.
At OIL2 (0.25 bytes/param), even a 4B-parameter model fits within
the memory constraints of a high-end smartphone.

\subsection{Differential Privacy Connections}

OIL training introduces an inherent source of differential privacy:
the quantisation map $Q$ destroys fine-grained gradient information,
acting as a built-in noise mechanism. The quantization noise has a
bounded $\ell_\infty$ norm equal to half the codebook step size,
providing a natural upper bound on the sensitivity of each
parameter update.

\begin{mylemma}[OIL Quantisation as DP Mechanism]{lem:oil_dp}
Let $\Delta$ denote the maximum codebook step size. Under the
Gaussian mechanism, OIL quantisation provides an effective noise
scale of $\sigma_{\text{eff}} \geq \Delta / 2$ per parameter
update, yielding an additive $(\epsilon, \delta)$-differential
privacy guarantee of
$\epsilon \leq \frac{\Delta \sqrt{2 \ln(1.25/\delta)}}{q \cdot \sigma_g}$
where $q$ is the sampling rate and $\sigma_g$ is the gradient
norm bound.
\end{mylemma}

This result means that OIL training provides a non-trivial privacy
guarantee ``for free,'' without the explicit noise injection
typically required in differentially private SGD.

\subsection{Communication Efficiency}

In federated learning, communication bandwidth is the primary
bottleneck. OIL4 reduces gradient communication by $8\times$
compared to FP32, while OIL2 achieves $16\times$ reduction. For
cross-device federated learning over cellular networks, this reduction
can transform an impractical protocol into a feasible one.

% -------------------------------------------------------------------------
\section{Hardware Acceleration}
\label{sec:hardware_accel}

While the current MYTHOS.cpp engine achieves high throughput on
general-purpose CPUs, specialised hardware could unlock an additional
order of magnitude in performance.

\subsection{FPGA Designs}

Field-programmable gate arrays are well-suited to OIL inference
because the codebook lookup operation maps directly to block RAM
(BRAM) content-addressable memory. An FPGA implementation of OIL8
inference would store each 256-entry codebook in a single BRAM
block and perform index-to-centroid lookup in a single clock cycle.
The pipelined architecture would achieve one multiply-accumulate
per cycle per BRAM pair, yielding a theoretical peak of
$\text{frequency} \times \text{BRAM pairs}$ operations per second.

\subsection{ASIC for Codebook Lookup}

A purpose-built ASIC for OIL inference would include:
\begin{itemize}
    \item \textbf{On-chip SRAM codebooks:} 256-entry FP32 or FP16
          codebooks with single-cycle access, sized for OIL8 and
          OIL4 respectively.
    \item \textbf{Nibble unpackers:} Hardwired logic to extract
          4-bit indices from packed byte arrays for OIL4, operating
          at one byte per cycle.
    \item \textbf{Popcount engines:} Massively parallel XOR-popcount
          arrays for Ternary GEMM, processing 1024 ternary weights
          per cycle.
    \item \textbf{Sparse activation units:} Hardware support for
          the MoE format heterogeneity described in
          Section~\ref{sec:architecture_codesign}.
\end{itemize}

\subsection{Memory Bandwidth Advantage}

The fundamental advantage of OIL on custom silicon is the reduction
in memory bandwidth demand. At OIL4, a 7B-parameter model requires
3.5\,GB of weight data, which fits in the HBM of a single
high-bandwidth memory stack. This eliminates the multi-chip
communication that dominates the cost of FP32 inference for large
models.

\begin{center}
\begin{tabular}{lrrr}
\toprule
\textbf{Format} & \textbf{7B Model Size} &
\textbf{HBM2 Stacks} & \textbf{Bandwidth Req.} \\
\midrule
FP32    & 28.0\,GB & 4 & 1.12\,TB/s \\
OIL8    &  7.0\,GB & 1 & 0.28\,TB/s \\
OIL4    &  3.5\,GB & 1 & 0.14\,TB/s \\
OIL2    &  1.75\,GB & 1 & 0.07\,TB/s \\
Ternary &  1.39\,GB & 1 & 0.06\,TB/s \\
\bottomrule
\end{tabular}
\end{center}

% -------------------------------------------------------------------------
\section{Comparison with Competing Approaches}
\label{sec:competing_approaches}

The quantised deep learning landscape includes several prominent
approaches. Understanding how OIL relates to each clarifies its
unique contributions.

\subsection{BitNet and 1-bit LLMs}

BitNet~\cite{wang2024} trains networks with ternary weights
$\{-1, 0, +1\}$ using an STE-based approach. OIL differs in three
fundamental ways. First, OIL provides a \emph{format family}
spanning 1.58 to 32\,BPW, allowing precision to be matched to the
task rather than fixed at 1-bit. Second, OIL's codebook-based
quantisation preserves gradient flow through the centroid values,
while BitNet relies on the STE which introduces gradient
approximation error. Third, OIL includes a formal information-theoretic
framework (the conservation law) that BitNet lacks.

\subsection{QLoRA and Low-Rank Adaptation}

QLoRA~\cite{dettmers2023} fine-tunes frozen quantised models using
low-rank adapters. This is a fundamentally different paradigm from
OIL: QLoRA treats quantisation as a compression step applied
\emph{after} or \emph{alongside} training, while OIL integrates
quantisation \emph{into} the training loop itself. The OIL gradient
dead zone provides implicit regularisation that QLoRA cannot replicate
because QLoRA's frozen weights never participate in the gradient
computation.

\subsection{DoRA and Weight-Decomposed Adaptation}

DoRA~\cite{liu2024} decomposes weights into magnitude and direction
components and adapts them separately. This decomposition could be
synergistic with OIL: the direction component, which is less
sensitive to quantisation, could use a lower BPW format, while the
magnitude component retains higher precision. This direction-magnitude
asymmetry is consistent with the observation that OIL's quantisation
error is dominated by radial (magnitude) perturbation rather than
angular (directional) perturbation.

\begin{mythm}[OIL vs.\ Post-Training Quantisation]{thm:oil_vs_ptq}
Let $\mathcal{L}_{\text{OIL}}(b)$ denote the loss achievable by
training from scratch in OIL format $b$, and let
$\mathcal{L}_{\text{PTQ}}(b)$ denote the loss after post-training
quantisation of an FP32-trained model to format $b$. Under the CID
assumption, there exists a crossover precision $b^*$ such that
$\mathcal{L}_{\text{OIL}}(b) < \mathcal{L}_{\text{PTQ}}(b)$ for
all $b < b^*$.
\end{mythm}

This theorem formalises the intuition that native quantised training
outperforms post-training quantisation at low bit-widths, because the
gradient dead zone acts as a stronger regulariser when the model is
trained from scratch in the quantised regime.

% -------------------------------------------------------------------------
\section{Open Problems and Conjectures}
\label{sec:open_problems}

We conclude with a summary of the most pressing open problems
organised by difficulty.

\begin{center}
\begin{tabular}{lp{5.8cm}c}
\toprule
\textbf{Problem} & \textbf{Description} & \textbf{Difficulty} \\
\midrule
Sub-1.5\,BPW training &
  Overcome gradient collapse below 1.5\,BPW without STE approximation. &
  Hard \\
Codebook co-optimisation &
  End-to-end differentiable codebook scheduling during training. &
  Medium \\
Formal convergence &
  Prove convergence of OIL training to stationary points. &
  Hard \\
Scaling law calibration &
  Empirically validate the modified Chinchilla laws of Section~\ref{sec:scaling_laws}. &
  Medium \\
MoE format allocation &
  Principled per-expert format selection based on routing statistics. &
  Medium \\
Federated OIL &
  Implement and evaluate OIL-based federated learning with DP guarantees. &
  Hard \\
ASIC tapeout &
  Fabricate a custom OIL inference chip and benchmark against GPU baselines. &
  Very Hard \\
\bottomrule
\end{tabular}
\end{center}

\subsection{Prioritised Roadmap}

The recommended research sequence is:

\begin{enumerate}
    \item \textbf{Short term (6 months):} Validate scaling laws
          empirically. Implement adaptive codebook scheduling.
          Benchmark MoE format heterogeneity on existing models.
    \item \textbf{Medium term (12 months):} Develop formal
          convergence proofs for OIL training. Implement federated
          OIL with differential privacy analysis. Produce FPGA
          prototype for OIL4 inference.
    \item \textbf{Long term (24+ months):} Explore sub-1.5\,BPW
          training via lottery-ticket-style sparse activation.
          Tapeout custom OIL ASIC. Integrate OIL training into
          production-scale (7B+) model pipelines.
\end{enumerate}

The research programme outlined here is ambitious but tractable. Each
direction builds on the theoretical foundations established in
Chapters~2--5 and the empirical evidence of Chapter~9, creating a
coherent path from the current proof-of-concept to a mature,
production-grade quantised training ecosystem.

\vfill
\begin{center}
\textit{``What we know is a drop, what we do not know is an ocean.''}\\[0.3em]
--- Isaac Newton
\end{center}
